\itemhead{习题3.2 · 第1题（教材p.85）}{XJ-B1-C03-HW-P085-Y0449}
\label{XJ-B1-C03-HW-P085-Y0449}
{\footnotesize 湘教版教材p.85 / PDF第92页 / 习题3.2 / 第1题}
\field{原题}{如图所示是函数 $f(x)$ 的图象，列出 $f(x)$ 的若干区间，说明它在各区间上的增减性，并指出该函数的最大、最小值点及最值。
\bankfig{C03-HW085-01}{0.68}}
\field{参考答案}{减区间 $[-5,-3],[-1,2],[3,5]$；增区间 $[-3,-1],[2,3]$。最大值 $1.5$ 在 $x=-1$；最小值 $-1$ 在 $x=2$。若用开区间列举也可，但最值点需单独写明。}
{\small \field{关联知识}{F04 函数三种表示法、图象与读数（教材帮p.85、86、93、94 / PDF94、95、102、103）；F05 单调性定义、证明与参数（教材帮p.97、98、100 / PDF106、107、109）；F06 最大值、最小值与值域（教材帮p.99、83、84 / PDF108、92、93）}
\field{方法与易错点}{把关键点与定义域一起标到图上，逐区间读交点和上下关系。 易错：将示意图预测当事实；把两曲线交点收进严格不等式；折线漏端点。}
\field{答案依据}{答案册相关题组：PDF第23、24页（册内页22、23）；按当前题库原题逐项重算，不按旧版题号机械对套。}}
